\documentclass[10pt,letterpaper]{article}
\usepackage[top=0.50in,bottom=0.52in,left=0.62in,right=0.62in,headheight=14pt,headsep=5pt,footskip=15pt]{geometry}
\usepackage{fontspec}
\setmainfont{Noto Serif}
\setsansfont{Noto Sans}
\usepackage[spanish,es-noshorthands]{babel}
\usepackage{microtype,ragged2e,setspace}
\usepackage{graphicx}
\graphicspath{{../compact-6p/figures/}}
\usepackage{booktabs,tabularx,array,multirow,longtable}
\usepackage[table]{xcolor}
\usepackage{tikz}
\usetikzlibrary{arrows.meta,positioning,calc,shapes.geometric,babel}
\usepackage{amsmath,amssymb,rotating}
\usepackage{fancyhdr,hyperref,enumitem,multicol}

\definecolor{autumn}{HTML}{E36A2E}
\definecolor{spicy}{HTML}{C8572A}
\definecolor{onyx}{HTML}{0C0C0C}
\definecolor{mustard}{HTML}{FCD757}
\definecolor{canary}{HTML}{8CB369}
\colorlet{orange}{autumn}
\colorlet{ink}{onyx}
\colorlet{muted}{onyx!58!white}
\colorlet{rule}{onyx!18!white}
\definecolor{soft}{HTML}{F5F7FA}
\definecolor{softblue}{HTML}{EAF2FD}
\definecolor{softgreen}{HTML}{EAF7EF}
\definecolor{softwarm}{HTML}{FDEFE8}
\definecolor{softplum}{HTML}{F2ECFB}
\definecolor{gamecol}{HTML}{2D6CC0}
\definecolor{controlcol}{HTML}{E36A2E}
\definecolor{searchcol}{HTML}{2E9B5F}
\definecolor{learncol}{HTML}{7B4DBA}
\definecolor{foundcol}{HTML}{6F7C8E}

\hypersetup{colorlinks=true,urlcolor=spicy,linkcolor=spicy,pdftitle={Marco teorico y estado del arte para el TFM MROB},pdfauthor={Jorge Luis Mayorga Taborda}}
\pagestyle{fancy}\fancyhf{}
\fancyhead[L]{\sffamily\fontsize{7.3}{8}\selectfont Jorge Luis Mayorga Taborda}
\fancyhead[R]{\sffamily\fontsize{7.3}{8}\selectfont TFM MROB · Marco teórico y estado del arte}
\fancyfoot[C]{\sffamily\fontsize{8.2}{8.7}\selectfont\thepage}
\renewcommand{\headrulewidth}{0.25pt}
\setlength{\parindent}{0pt}\setlength{\parskip}{2.2pt}
\setstretch{1.05}\color{ink}\emergencystretch=1.4em
\setlist[itemize]{leftmargin=4.3mm,itemsep=.8pt,topsep=.7pt,parsep=0pt}
\setlist[enumerate]{leftmargin=4.7mm,itemsep=.6pt,topsep=.7pt,parsep=0pt}

\newcommand{\sect}[1]{\par\vspace{0pt}{\sffamily\fontsize{14.5}{15.5}\selectfont\bfseries\color{orange}#1\par}\vspace{1.2pt}\textcolor{rule}{\rule{\linewidth}{0.35pt}}\vspace{3pt}}
\newcommand{\subh}[1]{\par\vspace{2.2pt}{\sffamily\fontsize{9.3}{10.2}\selectfont\bfseries\color{orange}#1\par}\vspace{.5pt}}
\newcommand{\figcap}[2]{\par\vspace{1.0pt}{\setstretch{1}\fontsize{8.1}{8.9}\selectfont\textbf{#1.}~\textit{#2}\par}\vspace{1.0pt}}
\newcommand{\sourcecap}[1]{\par{\setstretch{1}\fontsize{7.0}{7.7}\selectfont\itshape\centering Fuente: #1\par}}
\newcommand{\note}[1]{\par{\setstretch{1}\fontsize{7.0}{7.8}\selectfont\color{muted}#1\par}}
\newcommand{\citen}[1]{\textsuperscript{\,[#1]}}
\newcolumntype{L}[1]{>{\RaggedRight\arraybackslash}p{#1}}
\newcolumntype{C}[1]{>{\centering\arraybackslash}p{#1}}
\newcolumntype{Y}{>{\RaggedRight\arraybackslash}X}

\begin{document}
\setcounter{page}{5}

% ================= PAGE 5 =================
La comparación por periodos se apoya en el mismo corpus de mapeo: 3.014 identidades descubiertas, depuradas hasta 244 documentos analíticos. De ellos, 168 admitieron extracción estructural y 20 recibieron lectura cercana. Los 243 documentos fechados se agruparon en periodos de 35, 40, 79 y 89 registros. Las etiquetas proceden del título y el resumen, pueden coexistir y miden presencia temática, no rendimiento. La Figura 5 detalla el cambio metodológico y la Figura 6 comprueba si esa diversidad alcanza las interfaces de la misión.

\begin{center}
\includegraphics[width=0.985\linewidth]{fig06_method_evolution.pdf}
\end{center}
\figcap{Figura 5}{Evolución metodológica normalizada del corpus ampliado. Las categorías son multietiqueta y las cuotas se calculan dentro de cada periodo.}
\sourcecap{elaboración propia a partir del corpus analítico consolidado; observación parcial de 2026 y categorías coexistentes.}

\begin{center}
\includegraphics[width=0.985\linewidth]{fig07_interfaces_evidence.pdf}
\end{center}
\figcap{Figura 6}{Cobertura temática, intersecciones SP1--SP3 y madurez de validación observada en el corpus ampliado.}
\sourcecap{tablas de temas, combinaciones y validación, $n=244$; ``caso industrial'' puede describir simulación.}

\fontsize{8.35}{9.10}\selectfont
Leídas juntas, ambas figuras descartan una sustitución simple entre generaciones de métodos. Consenso/distribución pasa de 20,0\% a 24,7\%; subasta/mercado baja de 22,9\% a 13,5\%, pero continúa activa; aprendizaje sube de 2,9\% a 15,7\% y las metaheurísticas de 0\% a 11,2\%. La señal de CBF y seguridad alcanza 6,7\% en el periodo reciente. Selección, coordinación, control y seguridad se combinan con mayor diversidad, aunque siguen resolviendo tramos diferentes de la misión.

La cobertura por subproblemas conserva un desequilibrio mayor que cualquier cambio de cuota. SP1 aparece en 115 documentos, SP2 en 15 y SP3 en 20. Solo dos enlazan SP1--SP2 y diez SP1--SP3; ninguno conecta SP2--SP3 o cubre los tres. La evidencia también permanece segmentada: 99 documentos presentan simulación, 12 experimento físico y siete análisis formal. Solver, mercado, consenso y control ocupan etapas diferentes. Esta distribución desplaza el problema desde la elección de un método aislado hacia los contratos que permiten pasar de una etapa a la siguiente sin conservar una autoridad global permanente.
\normalsize

\newpage
% ================= PAGE 6 =================
La brecha identificada se traduce en tres contratos. SP1 entrega una coalición entera con roles definidos; SP2 certifica geometría, contacto, wrench y actuación; SP3 reserva recursos para la huella compuesta y realimenta el coste de congestión o fallo. Cada decisión debe superar el certificado de la capa siguiente. La ejecución es asíncrona, sin rondas globales y con periodo de muestreo explícito.

\figcap{Figura 7}{Contrato de interfaces entre formación de coaliciones, factibilidad física y coordinación del tráfico.}
\begin{center}
\begin{tikzpicture}[x=1mm,y=1mm,font=\sffamily]
\tikzset{
  stage/.style={rounded corners=1.5pt,draw=rule,line width=.45pt,minimum width=47mm,minimum height=23mm,fill=white,inner sep=0pt},
  stagehead/.style={rounded corners=1.1pt,minimum width=45.5mm,minimum height=5.4mm,text=white,align=center,font=\sffamily\fontsize{7.0}{7.4}\selectfont},
  stagebody/.style={align=center,font=\sffamily\fontsize{7.35}{8.25}\selectfont},
  state/.style={rounded corners=2.2pt,minimum width=39mm,minimum height=5.2mm,fill=white,align=center,font=\sffamily\fontsize{6.45}{6.8}\selectfont},
  checkpoint/.style={diamond,aspect=1.15,minimum size=4.8mm,inner sep=0pt,fill=orange,draw=white,line width=.45pt},
  flow/.style={-{Stealth[length=1.8mm]},draw=muted,line width=.65pt}
}
\node[stage,fill=softblue!62!white] (sp1) at (24,18) {};
\node[stage,fill=softwarm!65!white] (sp2) at (82.5,18) {};
\node[stage,fill=softgreen!70!white] (sp3) at (141,18) {};

\node[stagehead,fill=gamecol] at ([yshift=-3.25mm]sp1.north) {SP1 \enspace FORMACIÓN};
\node[stagehead,fill=controlcol] at ([yshift=-3.25mm]sp2.north) {SP2 \enspace TRANSPORTE};
\node[stagehead,fill=searchcol] at ([yshift=-3.25mm]sp3.north) {SP3 \enspace TRÁFICO};

\node[stagebody] at ([yshift=-.3mm]sp1.center) {\textbf{Decisión estratégica}\\capacidad, rol y compromiso};
\node[stagebody] at ([yshift=-.3mm]sp2.center) {\textbf{Factibilidad física}\\geometría, contacto, wrench y límites};
\node[stagebody] at ([yshift=-.3mm]sp3.center) {\textbf{Coordinación del tráfico}\\huella, reserva, prioridad y conflicto};

\node[state,draw=gamecol,text=gamecol] (formed) at ([yshift=4.15mm]sp1.south) {\textbf{Coalición completa}};
\node[state,draw=controlcol,text=controlcol] (feasible) at ([yshift=4.15mm]sp2.south) {\textbf{Margen físico válido}};
\node[state,draw=searchcol,text=searchcol] (completed) at ([yshift=4.15mm]sp3.south) {\textbf{Misión ejecutada}};

\node[checkpoint] (g12) at (53.25,18) {$\color{white}\checkmark$};
\node[checkpoint] (g23) at (111.75,18) {$\color{white}\checkmark$};
\draw[flow] (sp1.east) -- (g12.west);
\draw[flow] (g12.east) -- (sp2.west);
\draw[flow] (sp2.east) -- (g23.west);
\draw[flow] (g23.east) -- (sp3.west);

\coordinate (feedbackright) at ([yshift=-10.5mm]completed.south);
\coordinate (feedbackleft) at (feedbackright -| formed.south);
\draw[-{Stealth[length=1.8mm]},orange,line width=.65pt]
  (completed.south) -- (feedbackright)
  -- node[midway,below=1.1mm,fill=white,inner sep=1.2pt,font=\sffamily\fontsize{6.4}{6.8}\selectfont,text=spicy]
  {congestión, fallo o cambio de tarea: revisión local}
  (feedbackleft) -- (formed.south);
\end{tikzpicture}
\end{center}
\sourcecap{elaboración propia a partir de la discontinuidad observada entre SP1, SP2 y SP3; estados de ejecución según la cadena canónica del TFM.}

Esta cadena permite derivar el diseño experimental sin cambiar la unidad de comparación entre capas.
\rowcolors{2}{white}{soft}
\renewcommand{\arraystretch}{1.08}
\setlength{\tabcolsep}{2.6pt}
\fontsize{7.15}{7.75}\selectfont
\begin{tabularx}{\linewidth}{C{11mm} L{34mm} L{43mm} Y}
\toprule
\textbf{Capa} & \textbf{Mecanismo propuesto} & \textbf{Baseline/oráculo} & \textbf{Evidencia y métricas}\\
\midrule
SP1 & Juego local con unirse, esperar, cambiar rol o salir; payoff con compatibilidad, viaje, margen físico y congestión. & MILP/set partitioning como oráculo; CBBA/subasta distribuida; Húngaro solo si se demuestra reducción uno-a-uno. & Factibilidad, utilidad, gap, formación, churn, CPU, mensajes y bytes.\\
SP2 & Certificador de geometría, contacto, wrench y actuadores; pose/velocidad común y reparto de esfuerzo. & Leader--follower/estructura virtual; MPC/CBF; reparto central de wrench o impedancia. & Error de pose, residual de wrench, fuerza interna, contacto, saturación, energía y distancia mínima.\\
SP3 & Planificación de la huella carga--robots, negociación de corredores y retorno del coste de congestión a SP1. & A*/Dijkstra, priorizada, CBS/ECBS; ORCA/RVO o CBF-QP para seguridad reactiva. & Makespan, throughput, espera, conflictos, deadlock y replanteamientos.\\
Integral & Eventos locales y sustitución o disolución controlada ante fallos. & Oráculo end-to-end global, presentado como techo y no como comparación arquitectónica justa. & Entrega, recuperación, carga perdida, deadlines y dispersión entre semillas.\\
\bottomrule
\end{tabularx}
\normalsize

\fontsize{8.50}{9.25}\selectfont
Sobre este diseño se formulan cuatro hipótesis conectadas. H1 evalúa si el margen físico reduce coaliciones inviables con un coste acotado de utilidad. H2 contrasta si realimentar congestión disminuye makespan y bloqueos frente a asignar primero y enrutar después. H3 mide si la comunicación por eventos conserva factibilidad con menos bytes bajo pérdida y latencia acotadas. H4 compara la sustitución local con la disolución de la coalición cuando existe capacidad vecinal redundante.

Tres resultados publicados fijan escalas de referencia sin resolver la cadena completa. Liu y Shell demuestran optimalidad para matching uno-a-uno con $O(n^3\log n)$ bajo su estructura\citen{33}. Verma et al. escalan a 2.000 robots y 400 tareas, pero mediante una solución central sin física\citen{21}. Muhammed et al. validan dos AMR a 10 Hz, RMSE inferior a 4 cm y cerca de 70\% menos uso de red, sin reclutamiento ni tráfico\citen{26}. El contraste del TFM usará instancias idénticas, semillas pareadas, distribuciones, intervalos, fallos y ablaciones de margen físico, congestión y comunicación.

El alcance permanece condicionado por la revisión. WoS está incompleto en F01/F02, con 255 registros pendientes, y la API de arXiv no completó la cartera. Los porcentajes describen el corpus adquirido y solo 20 fuentes recibieron lectura cercana. La heterogeneidad de robots, cargas, métricas y escenarios impide un meta-análisis. En este corpus no se identificó una arquitectura que reúna coalición local heterogénea y variable, certificación física, tráfico de la huella completa y recuperación sin autoridad global permanente. La simulación del TFM aportará evidencia empírica, no una prueba de estabilidad, convergencia o escalabilidad global.

Las interfaces quedan así convertidas en objetos de validación: integridad de la coalición, factibilidad física y ejecución observable. Cada capa se comparará con su baseline pertinente y el oráculo global se declarará como techo de referencia, no como una comparación arquitectónica equivalente. La estructura bibliométrica aporta el último contraste: permite comprobar si la misma separación también aparece entre autores y comunidades temáticas.
\normalsize

\end{document}
